% LTeX: language=en-US

\section{Introduction}
\begin{frame}{Overview}
	\begin{wideitemize}
        \item Focus on models with \hi{occasionally binding constraints}
        \item Since Dynare 5: implementation of OccBin toolkit to solve models with occasionally binding constraints
        \item Feature only recently added\\
        $\rightarrow$ please report any bugs you may encounter
    \end{wideitemize}
\end{frame}

\begin{frame}{Occasionally binding constraints}
	\begin{wideitemize}
        \item Historically, economists relied on linearized models:
            \begin{wideitemize}
                \item Useful to gain intuition
                \item Easy to apply, especially for estimation (linear state space tools)
            \end{wideitemize}
        \item However: many interesting economic problems and policy questions feature nonlinearities such as \hi{occasionally binding constraints (OBCs)}:
        \begin{wideitemize}
            \item effective lower bound on nominal interest rates \parencite{ChristianoEichenbaumRebelo2011}
            \item Credit constraints \parencite{GuerrieriIacoviello2017}
            \item Downward nominal wage rigidities \parencite{SGU2016}
        \end{wideitemize}
        \item Not always sensible to assume OBC is always binding as in e.g.\ \textcite{Iacoviello2005} or \textcite{JermannQuadrini2012}
    \end{wideitemize}
\end{frame}


\begin{frame}{Occasionally Binding Constraints in Dynare}
	\begin{wideitemize}
        \item Dynare essentially offers three options to solve models with occasionally binding constraints:
        \begin{wideenumerate}
	    \item Nonlinear perfect foresight models: brute force (\texttt{max/min}-operators) and mixed complementarity problems (MCP)
	    \item Linearized stochastic models with surprise shocks: OccBin
        \item Stochastic nonlinear models: Extended Path
        \end{wideenumerate}
	    %\item Model estimations via Bayesian techniques or maximum likelihood, including models with OBCs
	\end{wideitemize}
\end{frame}


\begin{frame}{Crucial distinction: informational structure}

\begin{enumerate}
 \item deterministic/perfect foresight: all future shocks are known exactly
  \item stochastic simulations: only the shock distributions are known
\end{enumerate}

\begin{columns}[T]
\column{0.5\linewidth}
\begin{wideitemize}
    \item \hi{Deterministic:}
     \begin{itemize}
	        \item[+] Arbitrary precision
	        \item[+] Easily captures nonlinearities (such as OBCs)
	        \item[--] Assumption of perfect foresight
	        \item[--] How to simulate time series?
	    \end{itemize}
\end{wideitemize}
\column{0.5\linewidth}
\begin{wideitemize}
    \item \hi{Stochastic:}
	      \begin{itemize}
	        \item[+] Captures uncertainty
	        \item[+] Needed for estimation with state space form
            \item[--] Typically suffers from curse of dimensionality
	        \item[+] Perturbation methods can capture OBCs with additional tools (OccBin)
	        \item[--] Often uses 1st order approximations: certainty equivalence
	    \end{itemize}
\end{wideitemize}
\end{columns}
\end{frame}
